\problema{Word Ladder II}{126}{src/01-grafos/126-word-ladder-ii.cpp}{H}

Dadas \texttt{beginWord}, \texttt{endWord} y una lista \texttt{wordList},
cada palabra puede transformarse en otra cambiando exactamente una letra,
siempre que el resultado exista en \texttt{wordList}. A diferencia del
problema 127 (Word Ladder), donde sólo pedían la \emph{longitud} del camino
más corto, aquí hay que devolver \textbf{todas} las secuencias de
transformación más cortas de \texttt{beginWord} a \texttt{endWord}. Si
\texttt{endWord} es inalcanzable, se devuelve el vector vacío.

\paragraph{Intuición.} Primero resolvemos el 127 mentalmente: BFS por
niveles nos da la distancia mínima. La diferencia aquí es que, mientras
hacemos ese BFS, en vez de quedarnos sólo con la distancia, anotamos de
\emph{qué palabra viene cada palabra nueva} (su ``padre''). Como puede haber
varios padres válidos para una misma palabra en el mismo nivel, ese mapa de
padres es en realidad un grafo dirigido acíclico (DAG) por capas. Una vez
construido, basta con recorrerlo hacia atrás desde \texttt{endWord} hasta
\texttt{beginWord} para obtener, por construcción, únicamente caminos de
longitud mínima.

\begin{center}
\ttfamily
nivel 0:\ hit\\
nivel 1:\ hot\\
nivel 2:\ dot\quad lot\\
nivel 3:\ dog\quad log\\
nivel 4:\ cog
\end{center}

\noindent Con \texttt{wordList = ["hot","dot","dog","lot","log","cog"]} hay
dos caminos más cortos, ambos de longitud 5:
\texttt{hit-hot-dot-dog-cog} y \texttt{hit-hot-lot-log-cog}, porque tanto
\texttt{dog} como \texttt{log} tienen a \texttt{cog} como hija en el mismo
nivel.

\paragraph{Solución.} Dos fases:
\begin{enumerate}
  \item \textbf{BFS por niveles} desde \texttt{beginWord}. En cada nivel,
    para cada palabra probamos las 26 letras en cada posición; si el
    resultado está en el diccionario lo mandamos al siguiente nivel y
    registramos la palabra actual como uno de sus padres. Al terminar de
    procesar un nivel completo, borramos esas palabras del diccionario para
    que un nivel futuro no pueda \emph{regresar} a ellas (eso violaría la
    propiedad de camino más corto). Si \texttt{endWord} aparece en un nivel,
    terminamos de procesar ese nivel (puede tener varios padres) y detenemos
    el BFS.
  \item \textbf{DFS con backtracking} desde \texttt{endWord}, siguiendo
    todos los padres registrados hasta llegar a \texttt{beginWord},
    acumulando cada camino completo. Cada camino se arma en reversa y se
    invierte al final.
\end{enumerate}

\lstinputlisting{src/01-grafos/126-word-ladder-ii.cpp}

\complejidad{$O(N \cdot L^2 \cdot 26)$ para el BFS más el costo de reconstruir
todos los caminos (exponencial en el peor caso, inherente al problema)}
{$O(N \cdot L)$ para el diccionario y el mapa de padres}

\paragraph{Notas de entrevista (Amazon).} Este problema suele plantearse
como ``ya resolviste el 127, ahora dame todos los caminos''; la clave que
buscan es que reconozcas que un BFS normal con un solo padre por nodo pierde
caminos igualmente cortos, y que hay que permitir múltiples padres por nivel
antes de borrar las palabras visitadas. Errores típicos: borrar del
diccionario palabra por palabra en vez de por nivel completo (rompe caminos
válidos del mismo nivel), o mezclar BFS con DFS en una sola pasada (generar
todos los caminos con DFS puro, sin la poda de BFS, es exponencial e
inviable). Casos borde a mencionar: \texttt{endWord} no está en
\texttt{wordList} (vacío), \texttt{beginWord} no pertenece a
\texttt{wordList} (es válido, sólo debe empezar la cadena), y un único
camino más corto sin bifurcaciones.
